\documentclass[11pt,a4paper]{article}

\usepackage[margin=1.08in]{geometry}
\usepackage[T1]{fontenc}
\usepackage[utf8]{inputenc}
\usepackage{lmodern}
\usepackage{amsmath,amssymb,amsthm,mathtools,mathrsfs}
\usepackage{microtype}
\usepackage{enumitem}
\usepackage[colorlinks=true,linkcolor=blue,citecolor=blue,urlcolor=blue]{hyperref}

\newtheorem{theorem}{Theorem}[section]
\newtheorem{proposition}[theorem]{Proposition}
\newtheorem{lemma}[theorem]{Lemma}
\newtheorem{corollary}[theorem]{Corollary}
\theoremstyle{definition}
\newtheorem{definition}[theorem]{Definition}
\theoremstyle{remark}
\newtheorem{remark}[theorem]{Remark}

\DeclareMathOperator{\Fix}{Fix}
\DeclareMathOperator{\ord}{ord}
\DeclareMathOperator{\im}{im}
\DeclareMathOperator{\Rea}{Re}
\DeclareMathOperator{\Ima}{Im}

\newcommand{\ZZ}{\mathbb Z}
\newcommand{\RR}{\mathbb R}
\newcommand{\CC}{\mathbb C}
\newcommand{\BB}{\mathbf B}
\newcommand{\NN}{\mathbb N}
\newcommand{\cV}{\mathcal V}
\newcommand{\cP}{\mathcal P}
\newcommand{\cF}{\mathcal F}
\newcommand{\cS}{\mathcal S}
\newcommand{\cO}{\mathcal O}
\newcommand{\taup}{\tau}
\newcommand{\into}{\hookrightarrow}

\title{Spectral Packet Geometry of the Non-Trivial Zeta Zeros\
inside the Split-Zero Formalism}
\author{}
\date{May 2026}

\begin{document}
\maketitle

\begin{abstract}
Let
\[
S=G(\ZZ)=\ZZ\sqcup\{\taup\},\qquad A:=\{0,\taup\}\cong \BB,
\qquad F_1(s):=\zeta(s)-1.
\]
For every zero $\rho$ of the Riemann zeta function one has $F_1(\rho)=-1$, so the pointwise split-zero coefficient is the fixed Boolean pair $A$.  The non-trivial content is therefore not pointwise but geometric.  In this paper we determine the geometric image of the non-trivial zeros by combining the split-zero divisor formalism with the spectral constructions of Connes and Consani.

Let $\xi$ be the completed zeta function and let $D_\xi$ be its zero divisor.  We prove that the non-trivial zeros organize into Klein-four packets, that the Boolean multiplicity shadow descends to a packet sheaf on the orbit space, and that the affine spectral coordinate
\[
\Phi(s)=i\Bigl(s-\frac12\Bigr)
\]
identifies the non-trivial zero divisor with the spectral divisor cut out by $z\mapsto \zeta(\frac12-iz)$, equivalently with the common-zero locus of the Connes--Consani family $F(Ef)$.  Under this map the packet
\[
\{\rho,\bar\rho,1-\rho,1-\bar\rho\}
\]
becomes the centered spectral packet
\[
\{z,-z,\bar z,-\bar z\}.
\]
We show further that $\Phi$ is the square-root coordinate for the Connes Laplacian,
\[
-\rho(1-\rho)=-\Phi(\rho)^2-\frac14,
\]
and that on the critical locus the packet parameter determines the primitive $\zeta$-cycle length $2\pi/|\Phi(\rho)|$.  Finally, we prove the exact transformation law of the leading jet coefficient $\xi^{(m)}(\rho)/m!$ under the packet symmetries and obtain from it a well-defined positive packet weight.
\end{abstract}

\tableofcontents

\section{Introduction}

Let
\[
S=G(\ZZ)=\ZZ\sqcup\{\taup\},
\qquad
A:=\{0,\taup\}\subset S,
\qquad
F_1(s):=\zeta(s)-1.
\]
If $\rho$ is a zero of $\zeta$, then $F_1(\rho)=-1$, and multiplication by $-1$ fixes exactly the Boolean pair $A\cong\BB$.  Taken pointwise, every zeta zero therefore yields the same split-zero coefficient object.  The mathematical problem is not to distinguish zeros pointwise, but to determine the geometric object carried by the whole non-trivial zero divisor.

There is already a natural symmetry packet on the analytic side.  For the completed zeta function
\[
\xi(s):=\frac12 s(s-1)\pi^{-s/2}\Gamma\!\Bigl(\frac s2\Bigr)\zeta(s),
\]
complex conjugation and the functional involution $s\mapsto 1-s$ preserve the non-trivial zeros and their multiplicities.  Independently, Connes and Consani constructed two spectral realizations of the zeros of the Riemann zeta function: the first realization involves all non-trivial zeros and is expressed in terms of the Hochschild-trace map and the Laplacian $\Delta=H(1+H)$; the second involves the critical zeros and is expressed in terms of $\zeta$-cycles and sheaf cohomology on the Scaling Site.  The purpose of the present paper is to splice these two geometric structures into the split-zero formalism.

The main point is that the non-trivial zero divisor admits two compatible geometric descriptions.

\begin{enumerate}[label=(\arabic*)]
\item The quotient of the support by complex conjugation and $s\mapsto 1-s$ carries a canonical Boolean packet sheaf, whose stalk at a packet $O$ is the Boolean cube $A^{m_O}$ of the common multiplicity.
\item The affine map $\Phi(s)=i(s-\frac12)$ transports the same divisor to the spectral plane of Connes and Consani.  Under $\Phi$, the functional symmetry becomes $z\mapsto -z$ and complex conjugation becomes $z\mapsto -\bar z$.
\end{enumerate}

The outcome is a concrete answer to the question of how a non-trivial zero maps into the split-zero formalism.  After forgetting packet base and multiplicity one recovers the coefficient pair $A$; the full image is a \emph{weighted Boolean packet}: a spectral packet point in the Connes plane, a Boolean multiplicity cube attached to that packet, and on the critical locus a primitive circle length in the geometry of $\zeta$-cycles.

We keep the paper self-contained on the split-zero side.  On the Connes--Consani side we use, as inputs, their spectral realizations of all non-trivial zeros and of critical zeros \cite{ConnesConsaniCycles}.  For general background on sheaves and complex manifolds one may consult \cite{Lee}; for the Boolean and blueprint side of semiring geometry see \cite{Lorscheid}.  Nothing from these references beyond standard facts is used without explicit citation.

\section{The split-zero coefficient object}

\begin{definition}
Let
\[
S:=\ZZ\sqcup\{\taup\},\qquad \taup\notin \ZZ,
\]
with operations
\[
m\oplus n:=m+n,
\qquad
m\oplus\taup=\taup\oplus m:=m,
\qquad
\taup\oplus\taup:=\taup,
\]
and
\[
m\odot n:=mn,
\qquad
m\odot\taup=\taup\odot m:=\taup,
\qquad
\taup\odot\taup:=\taup.
\]
Thus $\taup$ is the additive identity and multiplicative absorber, while $1\in\ZZ\subset S$ is the multiplicative identity.
\end{definition}

\begin{definition}
Set
\[
A:=\{0,\taup\}\subset S.
\]
\end{definition}

\begin{proposition}\label{prop:boolean-pair}
The subset $A$ is a subsemiring of $S$, and the map
\[
\varphi:A\longrightarrow \BB=\{0,1\},
\qquad
\varphi(\taup)=0,
\qquad
\varphi(0)=1,
\]
is an isomorphism of semirings.
\end{proposition}

\begin{proof}
Inside $A$ one has
\[
0\oplus 0=0,
\qquad
0\oplus\taup=0,
\qquad
\taup\oplus\taup=\taup,
\]
and
\[
0\odot 0=0,
\qquad
0\odot\taup=\taup,
\qquad
\taup\odot\taup=\taup.
\]
These are exactly the Boolean addition and multiplication tables under the displayed identification.  Hence $A$ is a subsemiring and $\varphi$ is an isomorphism.
\end{proof}

\begin{proposition}\label{prop:units-fixed}
The unit group of $S$ is $S^\times=\{\pm 1\}$.  The element $-1$ is the unique unit of exact order $2$, and
\[
\Fix_S(-1)=A.
\]
\end{proposition}

\begin{proof}
Let $x\in S$ be a unit.  If $x=\taup$, then $x\odot y=\taup$ for all $y\in S$, so $x$ cannot be invertible.  Hence every unit lies in $\ZZ$.  The only units of $\ZZ$ are $\pm1$, so $S^\times=\{\pm1\}$.

Among these two units, only $-1$ has order $2$.  For the fixed-point set, first note that
\[
(-1)\odot\taup=\taup,
\]
so $\taup\in\Fix_S(-1)$.  If $n\in\ZZ$, then
\[
(-1)\odot n=n
\iff -n=n
\iff 2n=0
\iff n=0.
\]
Thus the only fixed supported point is $0$, and therefore
\[
\Fix_S(-1)=\{0,\taup\}=A.
\]
\end{proof}

\begin{corollary}\label{cor:F1-zero}
For every zero $\rho$ of the Riemann zeta function,
\[
F_1(\rho)=-1,
\qquad
\Fix_S(F_1(\rho))=A.
\]
\end{corollary}

\begin{proof}
If $\zeta(\rho)=0$, then
\[
F_1(\rho)=\zeta(\rho)-1=-1.
\]
Now apply Proposition \ref{prop:units-fixed}.
\end{proof}

\section{The non-trivial zero divisor and its packet space}

\subsection{The completed zeta function}

\begin{definition}
Define the completed zeta function
\[
\xi(s):=\frac12 s(s-1)\pi^{-s/2}\Gamma\!\Bigl(\frac s2\Bigr)\zeta(s).
\]
Let
\[
Z_\xi:=\{\rho\in\CC:\xi(\rho)=0\}
\]
be its zero set, and for $\rho\in Z_\xi$ let
\[
m_\rho:=\ord_\rho(\xi).
\]
The divisor of non-trivial zeta zeros is
\[
D_\xi:=\sum_{\rho\in Z_\xi} m_\rho[\rho].
\]
\end{definition}

\begin{lemma}\label{lem:trivial-simple}
For every integer $n\ge 1$, the point $-2n$ is a simple zero of $\zeta$.
\end{lemma}

\begin{proof}
Use the functional equation in the form
\[
\zeta(s)=2^s\pi^{s-1}\sin\!\Bigl(\frac{\pi s}{2}\Bigr)\Gamma(1-s)\zeta(1-s).
\]
At $s=-2n$, the factors $2^s$, $\pi^{s-1}$, $\Gamma(1-s)=\Gamma(1+2n)$, and $\zeta(1-s)=\zeta(2n+1)$ are holomorphic and nonzero.  The sine factor vanishes there, and its derivative equals
\[
\frac{d}{ds}\sin\!\Bigl(\frac{\pi s}{2}\Bigr)\Big|_{s=-2n}
=\frac{\pi}{2}\cos(-n\pi)
=\frac{\pi}{2}(-1)^n\neq 0.
\]
Hence the zero is simple.
\end{proof}

\begin{proposition}\label{prop:xi-entire}
The function $\xi$ is entire and satisfies
\[
\xi(s)=\xi(1-s).
\]
Its zeros are exactly the non-trivial zeros of $\zeta$, with the same multiplicities.
\end{proposition}

\begin{proof}
The functional equation for $\zeta$ gives the symmetry $\xi(s)=\xi(1-s)$ by direct substitution.

We check entireness.  Possible singularities can only come from the pole of $\zeta$ at $s=1$ and the poles of $\Gamma(s/2)$ at the nonpositive even integers.

At $s=1$, the factor $s-1$ cancels the simple pole of $\zeta$.
At $s=0$, the factor $s$ cancels the simple pole of $\Gamma(s/2)$.
At each negative even integer $-2n$ with $n\ge 1$, Lemma~\ref{lem:trivial-simple} shows that $\zeta$ has a simple zero and $\Gamma(s/2)$ has a simple pole, so these cancel.  No other singularities occur.  Hence $\xi$ is entire.

Now let $\rho$ be a zero of $\zeta$ that is not trivial.  The prefactor
\[
\frac12 s(s-1)\pi^{-s/2}\Gamma\!\Bigl(\frac s2\Bigr)
\]
is holomorphic and nonzero at $s=\rho$, so $\xi$ has a zero there of the same multiplicity.

Conversely, let $\rho$ be a zero of $\xi$.  The only zeros of the explicit prefactor are the factors $s$ and $s-1$, but the cancellations at $0$ and $1$ have already been accounted for in the proof of entireness, and these points are not zeros of $\xi$.  Therefore every zero of $\xi$ arises from a zero of $\zeta$.  The trivial zeros have been cancelled by the poles of $\Gamma(s/2)$, so the remaining zeros are exactly the non-trivial zeros of $\zeta$.
\end{proof}

\begin{lemma}\label{lem:real-zeros}
The real zeros of $\zeta$ are exactly the negative even integers.
\end{lemma}

\begin{proof}
If $s>1$ is real, then the Dirichlet series
\[
\zeta(s)=\sum_{n\ge 1} n^{-s}
\]
converges to a strictly positive real number, so $\zeta(s)\ne 0$.

If $0<s<1$, consider the Dirichlet eta function
\[
\eta(s):=\sum_{n\ge1} (-1)^{n-1} n^{-s}.
\]
For real $s\in(0,1)$, this is an alternating series with strictly decreasing positive terms tending to $0$, so $\eta(s)>0$.  Since
\[
\eta(s)=(1-2^{1-s})\zeta(s)
\]
and $1-2^{1-s}<0$ on $(0,1)$, one has $\zeta(s)<0$, in particular $\zeta(s)\ne 0$.

If $s<0$ is real, the functional equation gives
\[
\zeta(s)=2^s\pi^{s-1}\sin\!\Bigl(\frac{\pi s}{2}\Bigr)\Gamma(1-s)\zeta(1-s).
\]
Here $1-s>1$, so $\zeta(1-s)>0$, and $\Gamma(1-s)>0$ because $1-s$ is positive real.  The factors $2^s$ and $\pi^{s-1}$ are also nonzero.  Thus $\zeta(s)=0$ if and only if
\[
\sin\!\Bigl(\frac{\pi s}{2}\Bigr)=0,
\]
which for negative real $s$ occurs exactly at $s=-2,-4,-6,\dots$.
\end{proof}

\begin{corollary}\label{cor:nonreal}
Every zero of $\xi$ is non-real.
\end{corollary}

\begin{proof}
By Proposition \ref{prop:xi-entire}, the zeros of $\xi$ are exactly the non-trivial zeros of $\zeta$.  By Lemma \ref{lem:real-zeros}, the only real zeros of $\zeta$ are the trivial zeros.  Hence every zero of $\xi$ is non-real.
\end{proof}

\subsection{Klein-four packets}

\begin{definition}
Define maps on $\CC$ by
\[
c(s):=\bar s,
\qquad
w(s):=1-s.
\]
Let
\[
W:=\{\mathrm{id},c,w,cw\}.
\]
\end{definition}

\begin{lemma}\label{lem:conjugate-order}
Let $\Omega\subset\CC$ be a domain stable under complex conjugation, and let $f$ be holomorphic on $\Omega$ and satisfy
\[
f(\bar s)=\overline{f(s)}
\qquad (s\in \Omega).
\]
If $\rho\in\Omega$ is a zero of $f$, then
\[
\ord_{\bar\rho}(f)=\ord_\rho(f).
\]
\end{lemma}

\begin{proof}
Suppose
\[
f(s)=\sum_{n\ge m} a_n (s-\rho)^n
\]
near $\rho$, where $a_m\ne 0$.  Then near $\bar\rho$ one has
\[
f(s)=\overline{f(\bar s)}=
\overline{\sum_{n\ge m} a_n (\bar s-\rho)^n}
=
\sum_{n\ge m} \overline{a_n}(s-\bar\rho)^n.
\]
The first nonzero coefficient in this expansion is $\overline{a_m}\ne 0$, so the order of vanishing at $\bar\rho$ is again $m$.
\end{proof}

\begin{proposition}\label{prop:packet-symmetry}
If $\rho\in Z_\xi$, then
\[
\bar\rho,\qquad 1-\rho,\qquad 1-\bar\rho
\]
all lie in $Z_\xi$, and all four points occur with the same multiplicity.
\end{proposition}

\begin{proof}
The functional equation $\xi(1-s)=\xi(s)$ shows that $1-\rho$ is a zero whenever $\rho$ is a zero.  Since $w$ is biholomorphic with nonzero derivative, it preserves orders of vanishing, so
\[
\ord_{1-\rho}(\xi)=\ord_\rho(\xi).
\]

Next, on $\Rea(s)>1$ the Dirichlet series for $\zeta$ has real coefficients, so
\[
\zeta(\bar s)=\overline{\zeta(s)}
\]
on that half-plane; by analytic continuation this identity holds meromorphically on all of $\CC$.  The same is then true for $\xi$, because the factors $s(s-1)$, $\pi^{-s/2}$, and $\Gamma(s/2)$ are conjugation-compatible.  Thus
\[
\xi(\bar s)=\overline{\xi(s)}.
\]
By Lemma \ref{lem:conjugate-order}, $\bar\rho$ is a zero with the same multiplicity as $\rho$.  Applying both symmetries gives the same conclusion for $1-\bar\rho$.
\end{proof}

\begin{corollary}\label{cor:orbit-sizes}
Let $\rho\in Z_\xi$.
\begin{enumerate}[label=(\roman*)]
\item If $\Rea(\rho)=\frac12$, then the $W$-orbit of $\rho$ has cardinality $2$ and equals $\{\rho,\bar\rho\}$.
\item If $\Rea(\rho)\ne \frac12$, then the $W$-orbit of $\rho$ has cardinality $4$ and equals $\{\rho,\bar\rho,1-\rho,1-\bar\rho\}$.
\end{enumerate}
\end{corollary}

\begin{proof}
If $\Rea(\rho)=\frac12$, then $1-\rho=\bar\rho$, and likewise $1-\bar\rho=\rho$.  By Corollary \ref{cor:nonreal}, $\rho\ne\bar\rho$, so the orbit has exactly two points.

Now assume $\Rea(\rho)\ne\frac12$.  Then $\rho\ne1-\bar\rho$.  Again by Corollary \ref{cor:nonreal}$, \rho\ne\bar\rho$.  Also $\rho\ne1-\rho$, since that would force $\rho=1/2$.  Finally $\bar\rho\ne1-\rho$, for otherwise $\Rea(\rho)=1/2$.  Thus the four displayed points are pairwise distinct.
\end{proof}

\begin{definition}
Let
\[
Y_\xi:=Z_\xi/W
\]
be the orbit space of non-trivial zero packets, equipped with the quotient topology.  Because $Z_\xi$ is discrete, $Y_\xi$ is discrete as well.  For $O\in Y_\xi$, write $m_O$ for the common multiplicity of the points of the orbit.
\end{definition}

\subsection{The Boolean multiplicity shadow and packet descent}

\begin{lemma}\label{lem:product-sheaf}
Let $Z$ be a discrete topological space and let $(M_z)_{z\in Z}$ be a family of sets or $A$-semimodules.  For each open subset $U\subset Z$, define
\[
\cF_M(U):=\prod_{z\in U} M_z,
\]
with restriction maps given by coordinate projection.  Then $\cF_M$ is a sheaf.
\end{lemma}

\begin{proof}
Let $U\subset Z$ be open and let $(U_i)_{i\in I}$ be an open cover of $U$.  Because $Z$ is discrete, being equal on overlaps means equality of each coordinate indexed by a point of the overlap.

Suppose first that $x,y\in \cF_M(U)$ have equal restrictions to every $U_i$.  Let $z\in U$.  Choose $i$ with $z\in U_i$.  Equality on $U_i$ implies $x_z=y_z$.  Since this holds for every $z\in U$, one has $x=y$.

Now let $x_i\in \cF_M(U_i)$ be sections that agree on overlaps.  For each $z\in U$, choose $i$ with $z\in U_i$ and define $x_z$ to be the $z$-coordinate of $x_i$.  This is independent of the choice of $i$, because if $z\in U_i\cap U_j$, then the restrictions of $x_i$ and $x_j$ to the overlap are equal, so their $z$-coordinates coincide.  Thus $x=(x_z)_{z\in U}$ is a well-defined section of $\cF_M(U)$ restricting to each $x_i$.

Hence the sheaf axioms hold.
\end{proof}

\begin{definition}
The \emph{Boolean multiplicity shadow} of $D_\xi$ is the sheaf of $A$-semimodules on $Z_\xi$ defined by
\[
\cV_\xi(U):=\prod_{\rho\in U} A^{m_\rho}
\qquad (U\subset Z_\xi\ \text{open}).
\]
\end{definition}

\begin{definition}
Let $q_\xi:Z_\xi\to Y_\xi$ be the quotient map.  The \emph{Boolean packet sheaf} on $Y_\xi$ is
\[
\cP_\xi(U):=\prod_{O\in U} A^{m_O}
\qquad (U\subset Y_\xi\ \text{open}).
\]
\end{definition}

\begin{theorem}\label{thm:packet-descent-xi}
The group $W$ acts on the sheaf $\cV_\xi$ by reindexing the factors.  For every packet $O\in Y_\xi$ of cardinality $r_O$, the stalk of $q_{\xi,*}\cV_\xi$ at $O$ is canonically
\[
(q_{\xi,*}\cV_\xi)_O\cong (A^{m_O})^{r_O},
\]
and its $W$-fixed subsemimodule is the diagonal copy of $A^{m_O}$.  Consequently,
\[
(q_{\xi,*}\cV_\xi)^W\cong \cP_\xi
\]
as sheaves of $A$-semimodules on $Y_\xi$.
\end{theorem}

\begin{proof}
Because $Y_\xi$ is discrete, the stalk of $q_{\xi,*}\cV_\xi$ at a packet $O$ is the product of the stalks over the points in that orbit.  Since all points in the orbit have the same multiplicity by Proposition \ref{prop:packet-symmetry}, one obtains
\[
(q_{\xi,*}\cV_\xi)_O\cong \prod_{\rho\in O} A^{m_O}\cong (A^{m_O})^{r_O}.
\]

The action of $W$ on the orbit $O$ permutes the $r_O$ factors transitively.  An element of $(A^{m_O})^{r_O}$ is fixed by all such permutations if and only if all coordinates are equal.  The fixed subsemimodule is therefore the diagonal copy of $A^{m_O}$.

Since $Y_\xi$ is discrete, these diagonal stalks assemble to the sheaf whose value on an open set $U\subset Y_\xi$ is
\[
\prod_{O\in U} A^{m_O}=\cP_\xi(U).
\]
Hence $(q_{\xi,*}\cV_\xi)^W\cong \cP_\xi$.
\end{proof}

\section{The spectral plane and the Connes--Consani realization}

\subsection{The affine spectral coordinate}

\begin{definition}
Define the affine map
\[
\Phi:\CC\longrightarrow\CC,
\qquad
\Phi(s):=i\Bigl(s-\frac12\Bigr).
\]
Its inverse is
\[
\Phi^{-1}(z)=\frac12-iz.
\]
\end{definition}

\begin{proposition}\label{prop:Phi-symmetry}
For all $s\in\CC$ one has
\[
\Phi(1-s)=-\Phi(s),
\qquad
\Phi(\bar s)=-\overline{\Phi(s)}.
\]
Consequently $\Phi$ conjugates the action of $W$ on the $s$-plane to the action of the Klein four-group
\[
\widetilde W:=\{\mathrm{id},\sigma,\kappa,\sigma\kappa\}
\]
on the $z$-plane, where
\[
\sigma(z):=-z,
\qquad
\kappa(z):=-\bar z.
\]
In particular, if $\rho=\beta+i\gamma$, then
\[
\Phi(\rho)=-\gamma+i\Bigl(\beta-\frac12\Bigr),
\]
and the packet
\[
\{\rho,\bar\rho,1-\rho,1-\bar\rho\}
\]
is sent to
\[
\{z,-z,\bar z,-\bar z\},
\qquad z=\Phi(\rho).
\]
\end{proposition}

\begin{proof}
The identities are direct computations:
\[
\Phi(1-s)=i\Bigl(1-s-\frac12\Bigr)=i\Bigl(\frac12-s\Bigr)=-i\Bigl(s-\frac12\Bigr)=-\Phi(s),
\]
and
\[
\Phi(\bar s)=i\Bigl(\bar s-\frac12\Bigr)=-\overline{i\Bigl(s-\frac12\Bigr)}=-\overline{\Phi(s)}.
\]
Thus $w$ is conjugated to $\sigma$ and $c$ is conjugated to $\kappa$.  Applying these formulas to the four points of the $W$-orbit of $\rho$ yields the displayed spectral packet.
\end{proof}

\begin{definition}
Let
\[
Z_E:=\Phi(Z_\xi),
\qquad
D_E:=\Phi_*D_\xi=\sum_{z\in Z_E} m_z[z],
\\\qquad m_z:=m_{\Phi^{-1}(z)}.
\]
Let
\[
Y_E:=Z_E/\widetilde W
\]
be the spectral packet space.
\end{definition}

\begin{corollary}\label{cor:packet-space-iso}
The map $\Phi$ induces an isomorphism of discrete spaces
\[
\bar\Phi:Y_\xi\stackrel{\sim}{\longrightarrow} Y_E.
\]
\end{corollary}

\begin{proof}
By Proposition \ref{prop:Phi-symmetry}, $\Phi$ is equivariant for the isomorphism $W\cong\widetilde W$.  Since $\Phi$ is bijective on $\CC$, it restricts to a bijection $Z_\xi\to Z_E$ and therefore induces a bijection of quotient spaces.  Both spaces are discrete, so the induced bijection is a homeomorphism.
\end{proof}

\subsection{The Connes--Consani spectral divisor}

We now connect $D_E$ to the spectral realization of Connes and Consani.  Let $\cS(\RR)^{\mathrm{ev}}$ be the even Schwartz space on $\RR$, and let
\[
\cS(\RR)^{\mathrm{ev}}_1:=\left\{f\in\cS(\RR)^{\mathrm{ev}}: \int_{\RR} f(x)\,dx=0\right\}.
\]
For $f\in\cS(\RR)^{\mathrm{ev}}$, set
\[
(Ef)(u):=u^{1/2}\sum_{n\ge1} f(nu)
\qquad (u\in\RR_{>0}),
\]
and define its Mellin--Fourier transform by
\[
F(Ef)(z):=\int_{\RR_{>0}} (Ef)(u)u^{-iz}\,d^*u,
\qquad d^*u:=\frac{du}{u}.
\]

\begin{theorem}[Connes--Consani, \cite{ConnesConsaniCycles}]\label{thm:Connes-CC}
For every $f\in \cS(\RR)^{\mathrm{ev}}_1$ one has
\[
F(Ef)(z)=\zeta\Bigl(\frac12-iz\Bigr)\psi_f(z)
\]
on the half-plane $\Ima(z)>-\frac12$, where $\psi_f$ is holomorphic.  Consequently the common zero set of the functions $F(Ef)$, $f\in\cS(\RR)^{\mathrm{ev}}_1$, in the strip $|\Ima(z)|<\frac12$ is exactly the zero set of
\[
z\longmapsto \zeta\Bigl(\frac12-iz\Bigr).
\]
Moreover, the operator $\Delta=H(1+H)$ of \cite[Section~4]{ConnesConsaniCycles} has spectrum
\[
\left\{-\rho(1-\rho):\rho\in\CC\setminus\RR,\ \zeta(\rho)=0\right\}
\]
counted with multiplicity on the quotient described in \cite[Proposition~4.2]{ConnesConsaniCycles}.
\end{theorem}

\begin{proof}
This is exactly the content of equation~(2.2) and the common-zero statement in Section~2, together with Proposition~4.2 and its proof in Section~4 of \cite{ConnesConsaniCycles}.  The spectral values are written there as
\[
\Bigl(\rho-\frac12\Bigr)^2-\frac14,
\]
which is equal to $-\rho(1-\rho)$.
\end{proof}

\begin{theorem}\label{thm:bridge}
The affine map $\Phi$ identifies $Z_\xi$ with the Connes--Consani spectral support:
\[
Z_E=\left\{z\in\CC:\zeta\Bigl(\frac12-iz\Bigr)=0\right\}.
\]
Equivalently, $Z_E$ is the common zero set of the family $F(Ef)$, $f\in\cS(\RR)^{\mathrm{ev}}_1$, in the strip $|\Ima(z)|<\frac12$.

For every $\rho\in Z_\xi$ one has
\[
-\rho(1-\rho)=-\Phi(\rho)^2-\frac14.
\]
Thus $\Phi(\rho)$ is the square-root spectral coordinate for the Connes Laplacian.
\end{theorem}

\begin{proof}
Let $\rho\in Z_\xi$.  By Proposition \ref{prop:xi-entire}, $\rho$ is a non-trivial zero of $\zeta$.  Put $z=\Phi(\rho)$.  Since $\Phi^{-1}(z)=\frac12-iz$, one has
\[
\zeta(\rho)=0
\iff
\zeta\Bigl(\frac12-iz\Bigr)=0.
\]
Hence
\[
Z_E=\left\{z\in\CC:\zeta\Bigl(\frac12-iz\Bigr)=0\right\}.
\]
The strip condition is automatic, because if $\rho=\beta+i\gamma$ is non-trivial then $0<\beta<1$, and therefore
\[
\Ima(\Phi(\rho))=\beta-\frac12
\]
satisfies $|\Ima(\Phi(\rho))|<\frac12$.

The common-zero statement now follows from Theorem \ref{thm:Connes-CC}.

For the Laplacian relation, a direct calculation gives
\[
-\Phi(\rho)^2-\frac14
=-\Bigl(i\Bigl(\rho-\frac12\Bigr)\Bigr)^2-\frac14
=\Bigl(\rho-\frac12\Bigr)^2-\frac14
=-\rho(1-\rho).
\]
This is exactly the spectral parameter of $\Delta$ attached to $\rho$ by Theorem \ref{thm:Connes-CC}.
\end{proof}

\begin{definition}
The \emph{spectral Boolean packet sheaf} on $Y_E$ is
\[
\cP_E(U):=\prod_{P\in U} A^{m_P}
\qquad (U\subset Y_E\ \text{open}),
\]
where $m_P$ is the common multiplicity of the elements of the packet $P$.
\end{definition}

\begin{corollary}\label{cor:packet-transport}
Via the isomorphism $\bar\Phi:Y_\xi\to Y_E$, the packet sheaf $\cP_\xi$ transports canonically to $\cP_E$:
\[
\bar\Phi_*\cP_\xi\cong \cP_E.
\]
\end{corollary}

\begin{proof}
By definition of $D_E$, multiplicities are preserved by $\Phi$.  Hence for every open set $U\subset Y_E$ one has
\[
(\bar\Phi_*\cP_\xi)(U)=\cP_\xi(\bar\Phi^{-1}(U))
=\prod_{O\in \bar\Phi^{-1}(U)} A^{m_O}
=\prod_{P\in U} A^{m_P}
=\cP_E(U).
\]
These identifications commute with restriction maps.
\end{proof}

\section{Critical packets and $\zeta$-cycles}

\begin{proposition}\label{prop:critical-real}
Let $\rho\in Z_\xi$ and let $z=\Phi(\rho)$.  Then
\[
\Rea(\rho)=\frac12
\iff
z\in\RR.
\]
If $\rho=\frac12+i\gamma$, then
\[
\Phi(\rho)=-\gamma.
\]
Consequently the spectral packet of a critical zero is the two-point set
\[
\{-\gamma,\gamma\}.
\]
\end{proposition}

\begin{proof}
Write $\rho=\beta+i\gamma$.  Then Proposition \ref{prop:Phi-symmetry} gives
\[
\Phi(\rho)=-\gamma+i\Bigl(\beta-\frac12\Bigr).
\]
Thus $\Phi(\rho)$ is real if and only if $\beta=\frac12$.  In that case $\Phi(\rho)=-\gamma$.  The packet formula of Proposition \ref{prop:Phi-symmetry} then yields the two-point set $\{-\gamma,\gamma\}$.
\end{proof}

\begin{definition}
Let $Y_E^{\mathrm{crit}}\subset Y_E$ be the subset of critical packets, i.e. the packets contained in $\RR$.  For $P\in Y_E^{\mathrm{crit}}$, define
\[
\nu(P):=|z|
\]
for any $z\in P$.  This is well-defined because $P=\{z,-z\}$.

Define also the \emph{primitive cycle length}
\[
\Lambda(P):=\frac{2\pi}{\nu(P)}.
\]
\end{definition}

\begin{theorem}\label{thm:zeta-cycles}
Let $P\in Y_E^{\mathrm{crit}}$ and write $\nu(P)=s>0$.  Then every circle whose length is an integral multiple of
\[
\frac{2\pi}{s}
\]
is a $\zeta$-cycle in the sense of Connes and Consani.
\end{theorem}

\begin{proof}
Let $P=\{\pm s\}$ with $s>0$.  By Proposition \ref{prop:critical-real}, $P$ comes from a critical zero
\[
\rho=\frac12+is
\]
of $\zeta$.  Theorem~1.1(ii) of Connes and Consani, quoted in \cite{ConnesConsaniCycles}, states that if
\[
\zeta\Bigl(\frac12+is\Bigr)=0,
\]
then any circle whose length is an integral multiple of $2\pi/s$ is a $\zeta$-cycle.  Applying this with the displayed critical zero proves the assertion.
\end{proof}

\begin{corollary}\label{cor:critical-map}
The assignment
\[
\Lambda:Y_E^{\mathrm{crit}}\longrightarrow \RR_{>0},
\qquad
P\longmapsto \frac{2\pi}{\nu(P)},
\]
is a well-defined map from critical packets to primitive $\zeta$-cycle lengths.
\end{corollary}

\begin{proof}
This is immediate from the definition of $\nu(P)$ and Theorem \ref{thm:zeta-cycles}.
\end{proof}

\begin{remark}
The critical part of the packet map is therefore geometric in the strict Connes--Consani sense: besides the Boolean coefficient object, a critical zero packet determines a primitive circle length for the $\zeta$-cycles on the Scaling Site, together with all of its covering multiples.
\end{remark}

\section{Packet jet weights}

\begin{definition}
For $\rho\in Z_\xi$, let $m=m_\rho$.  Define the leading jet coefficient
\[
\lambda_\rho:=\frac{\xi^{(m)}(\rho)}{m!}.
\]
Since $m=\ord_\rho(\xi)$, one has $\lambda_\rho\ne 0$.
\end{definition}

\begin{lemma}\label{lem:local-expansion}
Let $\rho\in Z_\xi$ and $m=m_\rho$.  Then near $\rho$ one has the local expansion
\[
\xi(s)=\lambda_\rho (s-\rho)^m + O\bigl((s-\rho)^{m+1}\bigr).
\]
\end{lemma}

\begin{proof}
This is the Taylor expansion of $\xi$ at $\rho$.  By definition of the vanishing order, the derivatives of orders $0,1,\dots,m-1$ vanish at $\rho$, while the derivative of order $m$ does not.  Hence the first nonzero term is exactly
\[
\frac{\xi^{(m)}(\rho)}{m!}(s-\rho)^m=\lambda_\rho (s-\rho)^m.
\]
\end{proof}

\begin{theorem}\label{thm:lambda-transform}
Let $\rho\in Z_\xi$ and let $m=m_\rho$.  Then the leading jet coefficients satisfy
\[
\lambda_{1-\rho}=(-1)^m\lambda_\rho,
\qquad
\lambda_{\bar\rho}=\overline{\lambda_\rho},
\qquad
\lambda_{1-\bar\rho}=(-1)^m\overline{\lambda_\rho}.
\]
Consequently
\[
|\lambda_{1-\rho}|=|\lambda_{\rho}|,
\qquad
|\lambda_{\bar\rho}|=|\lambda_\rho|,
\qquad
|\lambda_{1-\bar\rho}|=|\lambda_\rho|.
\]
\end{theorem}

\begin{proof}
Write $m=m_\rho$.

By Lemma \ref{lem:local-expansion}, near $\rho$ one has
\[
\xi(\rho+t)=\lambda_\rho t^m + O(t^{m+1}).
\]
Using the functional equation $\xi(1-s)=\xi(s)$, one gets near $1-\rho$
\[
\xi(1-\rho+u)=\xi(\rho-u)=\lambda_\rho (-u)^m + O(u^{m+1})
=(-1)^m\lambda_\rho u^m + O(u^{m+1}).
\]
Therefore
\[
\lambda_{1-\rho}=(-1)^m\lambda_\rho.
\]

For conjugation, write near $\rho$
\[
\xi(s)=\sum_{n\ge m} a_n (s-\rho)^n,
\qquad a_m=\lambda_\rho.
\]
Since $\xi(\bar s)=\overline{\xi(s)}$, one has near $\bar\rho$
\[
\xi(s)=\overline{\xi(\bar s)}=
\sum_{n\ge m} \overline{a_n}(s-\bar\rho)^n.
\]
Thus the leading coefficient at $\bar\rho$ is $\overline{a_m}=\overline{\lambda_\rho}$, i.e.
\[
\lambda_{\bar\rho}=\overline{\lambda_\rho}.
\]
Applying both identities yields
\[
\lambda_{1-\bar\rho}=(-1)^m\lambda_{\bar\rho}=(-1)^m\overline{\lambda_\rho}.
\]
The absolute-value statements are immediate.
\end{proof}

\begin{definition}
For a packet $P\in Y_E$, choose any $\rho\in Z_\xi$ mapping to $P$ and set
\[
\omega(P):=|\lambda_\rho|=\left|\frac{\xi^{(m_P)}(\rho)}{m_P!}\right|.
\]
We call $\omega(P)$ the \emph{packet jet weight}.
\end{definition}

\begin{corollary}\label{cor:weight-well-defined}
The packet jet weight is well-defined.  Thus
\[
\omega:Y_E\longrightarrow \RR_{>0}
\]
is a well-defined positive function on the spectral packet space.
\end{corollary}

\begin{proof}
Any two representatives of a packet differ by the action of $W$.  Theorem \ref{thm:lambda-transform} shows that the absolute value of the leading jet coefficient is constant on each packet.  Therefore $\omega$ is independent of the chosen representative.
\end{proof}

\section{The geometric image of a non-trivial zero}

We now assemble the preceding constructions into a single theorem.

\begin{theorem}\label{thm:main}
Let $\rho\in Z_\xi$ be a non-trivial zero of the Riemann zeta function, and let
\[
P_\rho:=[\Phi(\rho)]\in Y_E
\]
be its spectral packet.
Then the following statements hold.

\begin{enumerate}[label=(\arabic*)]
\item The split-zero coefficient attached to $\rho$ is
\[
F_1(\rho)=-1,
\qquad
\Fix_S(F_1(\rho))=A=\{0,\taup\}\cong\BB.
\]

\item The packet base point is represented by
\[
\Phi(\rho)=i\Bigl(\rho-\frac12\Bigr),
\]
and the full analytic packet
\[
\{\rho,\bar\rho,1-\rho,1-\bar\rho\}
\]
becomes the centered spectral packet
\[
\{z,-z,\bar z,-\bar z\},
\qquad z=\Phi(\rho).
\]

\item The Boolean multiplicity shadow descends to the spectral packet sheaf $\cP_E$, and the stalk at $P_\rho$ is
\[
(\cP_E)_{P_\rho}\cong A^{m_\rho}.
\]
Thus the multiplicity of $\rho$ is recorded as the Boolean cube dimension of its packet stalk.

\item The Connes Laplacian sees the same zero through the spectral identity
\[
-\rho(1-\rho)=-\Phi(\rho)^2-\frac14.
\]

\item If $\rho=\frac12+i\gamma$ lies on the critical line, then
\[
P_\rho=\{\pm\gamma\}
\]
up to the sign convention imposed by $\Phi$, and the primitive $\zeta$-cycle length attached to the packet is
\[
\Lambda(P_\rho)=\frac{2\pi}{|\gamma|}.
\]

\item The packet jet weight
\[
\omega(P_\rho)=\left|\frac{\xi^{(m_\rho)}(\rho)}{m_\rho!}\right|
\]
is well-defined.
\end{enumerate}

In particular, the geometric image of a non-trivial zero in the split-zero formalism is the weighted Boolean packet
\[
\bigl(P_\rho,(\cP_E)_{P_\rho},\omega(P_\rho)\bigr),
\]
and on the critical locus this packet carries in addition the primitive $\zeta$-cycle length $\Lambda(P_\rho)$.
\end{theorem}

\begin{proof}
Item (1) is Corollary \ref{cor:F1-zero}.  Item (2) is Proposition \ref{prop:Phi-symmetry}.  Item (3) follows from Corollary \ref{cor:packet-transport} and the definition of $\cP_E$.  Item (4) is the second statement of Theorem \ref{thm:bridge}.  Item (5) is Proposition \ref{prop:critical-real} together with Corollary \ref{cor:critical-map}.  Item (6) is Corollary \ref{cor:weight-well-defined}.
\end{proof}

\begin{remark}
The theorem gives an exact answer to the packet-geometric question.  The constant coefficient object $A$ is the reduced coefficient remembered after forgetting packet base and multiplicity.  The full image is a packet point in the Connes spectral plane, equipped with a Boolean multiplicity cube and a canonical positive local weight.  On the critical locus, the same packet point determines the covering-stable family of $\zeta$-cycles on the Scaling Site.
\end{remark}

\begin{thebibliography}{99}

\bibitem{ConnesConsaniCycles}
A.~Connes and C.~Consani,
\emph{Hochschild homology, trace map and $\zeta$-cycles},
Proc. Sympos. Pure Math. \textbf{105} (2023), 83--104.

\bibitem{ConnesConsaniScaling}
A.~Connes and C.~Consani,
\emph{Geometry of the scaling site},
Selecta Math. (N.S.) \textbf{23} (2017), no.~3, 1803--1850.

\bibitem{ConnesTraceFormula}
A.~Connes,
\emph{Trace formula in noncommutative geometry and the zeros of the Riemann zeta function},
Selecta Math. (N.S.) \textbf{5} (1999), no.~1, 29--106.

\bibitem{Lee}
J.~M.~Lee,
\emph{Introduction to Complex Manifolds},
Graduate Studies in Mathematics, vol.~244, American Mathematical Society, Providence, RI, 2024.

\bibitem{Lorscheid}
O.~Lorscheid,
\emph{The geometry of blueprints. Part I: Algebraic background and scheme theory},
Adv. Math. \textbf{229} (2012), no.~3, 1804--1846.

\bibitem{Titchmarsh}
E.~C.~Titchmarsh,
\emph{The Theory of the Riemann Zeta-Function},
2nd ed., revised by D.~R.~Heath-Brown, Oxford University Press, New York, 1986.

\end{thebibliography}

\end{document}
